\section{Feasibility: Alchemy, Flight, or Perpetual Motion}
\label{sec:feasibility}

The templates introduced in Section~\ref{sec:intro} can now be adjudicated
against evidence. The strongest argument that decoding animal signals is
flight rather than alchemy is that it has already been done, repeatedly, in
other species. Vervet monkeys produce predator-specific alarm calls that
elicit appropriate escape responses, and habituation--dishabituation
experiments show the calls function as referential
labels~\cite{seyfarth1980vervet,cheney1990monkeys}. Prairie dogs encode
predator type, size, and color~\cite{slobodchikoff2009prairie}, chicken alarm
calls carry graded referential content~\cite{evans1993alarm}, the honeybee
waggle dance communicates distance and direction~\cite{vonfrisch1967dance},
bottlenose dolphins use signature whistles as individual
labels~\cite{janik1998whistles,king2013labels}, and chimpanzees inform
conspecifics of danger in a goal-directed way~\cite{crockford2012inform}.
Japanese tits go further, combining calls into compositional
sequences~\cite{suzuki2016compositional}, and formal-semantics analyses find
context-dependent, rule-governed meaning in animal
calls~\cite{schlenker2016formal}. Referential and compositional signaling
exists in nature, so decoding it is an engineering problem, not a category
error. The open question is whether dogs sit in this referential class; the
canine evidence places them closer to the affective end.

The dog, meanwhile, is a willing and readable partner. Domestication
reshaped its social cognition for life with
humans~\cite{hare2002domestication,frantz2016genomic}: dogs follow human
pointing and gaze better than chimpanzees~\cite{soproni2001comprehension},
look back to humans at unsolvable tasks~\cite{miklosi2003lookback}, show
reward-system activation to praise in fMRI~\cite{berns2012fmri}, possess
voice-sensitive brain regions homologous to
humans'~\cite{andics2014voice}, and lateralize tail-wagging by emotional
valence~\cite{quaranta2007tail}. A few dogs learn hundreds of object
names~\cite{kaminski2004rico,pilley2011chaser}, and soundboard-trained dogs
produce non-random, non-imitative two-button combinations and respond to
human button presses above
chance~\cite{bastos2024plos,bastos2024scirep,bastos2023interact}. The dog is
not a black box: it perceives human communicative intent, signals its own
states readably, and in controlled settings exchanges referential tokens.
Affect-and-intent decoding has empirical ground to stand on; it is flight
with a runway already built.

The pessimistic half of the verdict concerns Rung~5. The comparative
evidence suggests that the recursive, generative syntax human translation
presupposes, the ``faculty of language in the narrow sense,'' may be
uniquely human~\cite{hauser2002faculty,chomsky1957syntactic,hauser1996evolution}.
On this view the language instinct is a species-typical
adaptation~\cite{pinker1994instinct} rooted in shared
intentionality~\cite{tomasello2008origins}, and the cultural-evolution
literature shows how iterated learning and cumulative culture shape
linguistic structure into something no individual animal
invents~\cite{kirby2008cumulative,smith2003iterated}; even the statistical
universals of human language are signatures of a generative
system~\cite{zipf1949human}. If a dog's repertoire lacks this structure,
``translating its sentences'' is alchemy-like: not hard engineering but a
target that may not exist. The modality matters: this is not a
perpetual-motion proof. No conservation law forbids richer animal
signaling, and the compositionality results above keep the empirical
question open. Rung~5 is possibly ill-posed, not provably impossible.

The layered verdict, then: Rungs~1--2 are achieved flight; Rung~3 is flight
with a foothold; Rung~4 is flight in some species and an open question for
dogs; Rung~5 is alchemy-like. Nothing here is perpetual motion, and the
existence proofs guarantee that something real is achievable. The practical
consequence is to pour effort into the grounded, falsifiable lower rungs and
treat the upper rung as a hypothesis about the animal rather than a product
roadmap (Table~\ref{tab:proofs}).

\begin{table}[htbp]
\centering
\caption{Existence proofs and the feasibility template each supports.
Nothing here is ``perpetual motion'' (provably impossible).}
\label{tab:proofs}
\footnotesize
\setlength{\tabcolsep}{4pt}
\begin{tabular}{@{}p{3.4cm}p{8.0cm}p{2.2cm}@{}}
\toprule
\textbf{Claim} & \textbf{Existence proof} & \textbf{Template} \\
\midrule
Affect decodable (R2) & humans classify barks/growls~\cite{pongracz2005barks,farago2017growls}; tail lateralization~\cite{quaranta2007tail}; voice regions~\cite{andics2014voice} & flight (achieved) \\
Intent inferable (R3) & pointing~\cite{soproni2001comprehension}; look-back~\cite{miklosi2003lookback}; soundboard combos~\cite{bastos2024scirep} & flight (foothold) \\
Reference in animals (R4) & vervet \& prairie alarms~\cite{seyfarth1980vervet,slobodchikoff2009prairie}; dolphin labels~\cite{janik1998whistles}; bee dance~\cite{vonfrisch1967dance} & flight (species-dependent) \\
Compositionality (R4) & tit compositional calls~\cite{suzuki2016compositional}; formal semantics~\cite{schlenker2016formal} & flight (rare) \\
Dogs read human speech & fMRI reward~\cite{berns2012fmri}; word learning~\cite{kaminski2004rico,pilley2011chaser} & flight (achieved) \\
Open translation (R5) & no recursive syntax in dogs; narrow faculty possibly unique~\cite{hauser2002faculty,chomsky1957syntactic} & alchemy-like \\
Provably impossible? & no conservation law forbids signaling & not perpetual motion \\
\bottomrule
\end{tabular}
\end{table}
